\documentclass[11pt]{article}
\usepackage[T1]{fontenc}
\usepackage{lmodern,amsmath,amssymb,amsthm}
\usepackage[margin=1in]{geometry}
\usepackage[hidelinks]{hyperref}
\newcommand{\Q}{\mathbf Q}
\newcommand{\Z}{\mathbf Z}
\newcommand{\F}{\mathbf F}
\newcommand{\Sha}{\mathrm{Sha}}
\newcommand{\Sel}{\mathrm{Sel}}
\newcommand{\GL}{\mathrm{GL}}
\newtheorem{theorem}{Theorem}
\title{Splitting Densities of Mordell--Weil and Selmer Classes}
\author{\small\href{mailto:admin@cosmosgeek.org}{admin@cosmosgeek.org}
\quad\href{mailto:b321014@iiit-bh.ac.in}{b321014@iiit-bh.ac.in}}
\date{Working draft}
\begin{document}\maketitle
\begin{abstract}
For elliptic torsion representations satisfying two explicit image
hypotheses, we compute the joint evaluation field of any finite module of
cohomology classes, including nonfree modules. The quotient of the full
Selmer and full point splitting densities at level $p^n$ determines the
order of $\Sha[p^n]$. As $n$ varies, these ratios recover its finite
elementary divisors and divisible corank. For full $p$-adic image we also
evaluate the density of simultaneous divisibility of $r$ marked independent
points at all $p$-power levels and compute a nonconstant character
coefficient of the finite-level density. An explicit family illustrates
the difference between fields for the marked points and the full group.
\end{abstract}

\section{The evaluation field}
Let $p$ be odd, $n\ge1$, $R=\Z/p^n$, and let $V=E[p^n]$ for an
elliptic curve over $\Q$. Write $K=\Q(V)$, $G=\operatorname{Gal}(K/\Q)$,
and regard $V$ as a free $R$-module of rank two. Assume that $G$
contains a scalar $zI$ with $z-1\in R^\times$ and that the $R$-algebra
generated by $G$ is $\operatorname{End}_R(V)$. These conditions hold
when $G=\GL_2(R)$: use $-I$ and take differences of elementary matrices
and the identity to obtain the off-diagonal matrix units; their products
give the diagonal units.

Let $A$ be a finite $R$-submodule of $H^1(\Q,V)$. Restriction of each
class to $\operatorname{Gal}(\overline{\Q}/K)$ is a continuous
homomorphism into $V$, independent of the cocycle representative.
Let $F_A$ be the joint fixed field of their kernels inside
$\overline{\Q}$. It is a finite Galois extension of $\Q$ containing $K$.

\begin{theorem}
Evaluation induces a canonical isomorphism
\[
 \operatorname{Gal}(F_A/K)\simeq\operatorname{Hom}_R(A,V).
\]
In particular $[F_A:\Q]=|G|\,|A|^2$.
\end{theorem}
\begin{proof}
On any cocycle $c$ of $G$, centrality gives
$(z-1)c(g)=(g-1)c(zI)$. Thus $H^1(G,V)=0$. By inflation--restriction,
a class vanishing on the absolute Galois group of $K$ vanishes globally.
Write $A^*=\operatorname{Hom}_R(A,R)$ and choose a basis of $V$.
The evaluation image $W$ lies in $(A^*)^2$. It is an additive subgroup,
hence an $R$-submodule. Conjugation acts on the output coordinates
through $G$ (the restrictions are restrictions of global cocycles).
This also proves the normality defining $F_A/\Q$. By the matrix-algebra
hypothesis $W$ is stable under all matrix units, so $W=U\oplus U$ for
some $U\subseteq A^*$.

The pairing $A\times A^*\to R$ is perfect as a pairing of finite
abelian groups: compose with $x\mapsto x/p^n\bmod\Z$ and use the
cyclic decomposition of $A$. If $U$ were proper its annihilator
in $A$ would be nonzero. Every element there would restrict to zero
on the absolute Galois group of $K$, contrary to injectivity.
Hence $W=(A^*)^2$. In particular $|A^*|=|A|$, including when $A$
has cyclic factors of smaller exponent.
\end{proof}

\section{Recovery of the Tate--Shafarevich group}
Let $M,S$ denote the images of $E(\Q)/p^nE(\Q)$ and
$\Sel_{p^n}(E/\Q)$ in $H^1(\Q,V)$. Then $M\subseteq S$ and
$S/M=\Sha(E/\Q)[p^n]$. For $B=M,S$, write $\delta_{B,n}$ for the
natural prime density of total splitting in $F_B$. Chebotarev gives
$\delta_{B,n}=1/[F_B:\Q]$. The theorem yields
\[
 [F_S:F_M]=|\Sha(E/\Q)[p^n]|^2,\qquad
 \frac{\delta_{M,n}}{\delta_{S,n}}=|\Sha(E/\Q)[p^n]|^2 .
\]
More precisely, over $F_M$ the evaluation group is
$\operatorname{Hom}_R(S/M,V)$. For $c\in S$ of order $p^e$ modulo
$M$, evaluation at its coset maps this group onto $V[p^e]$.
Indeed the character of a cyclic subgroup extends to its ambient finite
abelian group by divisibility of $\Q/\Z$, and the exponent bound keeps
all its values in $R$. Apply this argument coordinatewise. Among
primes splitting in $F_M$, the proportion with nonzero evaluation
at $c$ is $1-p^{-2e}$, since this test is invariant under conjugacy.
At unramified good primes away from $p$ this is the test for nonzero
localization of $c$ even though every global point is locally divisible
by $p^n$.

Suppose the image hypotheses hold for every $n$. Put
$h_0=0$ and $h_n=\tfrac12\log_p(\delta_{M,n}/\delta_{S,n})$.
The usual cofinitely generated structure of the $p$-primary
Tate--Shafarevich group, obtained as a quotient of the cofinitely
generated $p^\infty$-Selmer group, gives
\[
 \Sha[p^\infty]\simeq(\Q_p/\Z_p)^t\oplus
 \bigoplus_j\Z/p^{e_j},\qquad
 h_n=nt+\sum_j\min(n,e_j).
\]
Writing $b_n=h_n-h_{n-1}$, one gets $t=\lim b_n$ and
$\#\{j:e_j=n\}=b_n-b_{n+1}$. This reconstruction permits an infinite
divisible part. The data use the full rational point module and the full
Selmer module, so they need not be effectively known for a given curve.

\section{Divisibility with the linear part unrestricted}
Assume that $E$ has full $p$-adic image and that $P_1,\ldots,P_r$ are
independent modulo $p$. There is no rational $p$-torsion, so their
span at level $p^n$ is $R^r$ (their $\Z_p$-span is a direct summand).
Evaluation gives full translations $V^r$ over $K$. For fixed linear
Frobenius $g$, the chance over this translation fibre that
all the $P_i$ are locally divisible by $p^n$ is
$|\operatorname{coker}(I-g:V\to V)|^{-r}$. Indeed a chosen division
point can be translated to a Frobenius-fixed point precisely when
its displacement belongs to $(I-g)V$. The affine action also explains
why division statistics for independent Selmer classes alone cannot
distinguish their rational origin.

At residual level put $g_p=p(p-1)^2(p+1)$ and let $N_j$ be the number
of invertible matrices with $j$-dimensional fixed space. Then
\[
 N_2=1,\qquad N_1=(p+1)(p^2-p-1),\qquad N_0=g_p-N_1-1.
\]
For $N_1$ the distinct eigenvalues $1,\lambda\ne1$ give $p(p+1)$
matrices per choice of $\lambda$, and the nonidentity unipotents give
$p^2-1$. Thus the density at level one is
$(N_0+N_1p^{-r}+p^{-2r})/g_p$. If one also requires the reducing prime
$q$ to satisfy $q\equiv1\bmod p$ and $a_q(E)\equiv2\bmod p$, the
density among all primes is $((p^2-1)p^{-r}+p^{-2r})/g_p$.

For the intersection of the divisibility conditions over all $n$, the
natural density is the probability integral
\[
 D_\infty(r)=\int_{\GL_2(\Z_p)}|\det(I-g)|_p^r\,dg.
\]
Indeed failure beyond level $n$, after success there, requires
$p^{n+1}\mid \#E(\F_q)$. Chebotarev bounds the upper density of
this tail by the Haar measure of $\{\det(I-g)\in p^{n+1}\Z_p\}$.
That measure tends to zero: the zero set of this nonzero polynomial
has measure zero by Fubini. This justifies the passage from finite
Chebotarev statements to natural density.

Put $C_r=(1-p^{-1})/(1-p^{-1-r})$. In corank one modulo $p$,
a cofactor of $I-g$ is a unit; the lifted determinant is uniform on
$p\Z_p$, so its moment is $p^{-r}C_r$. In corank two write
$g=I+pB$, with $B$ Haar on all $2$-by-$2$ matrices over $\Z_p$.
The first column of $B$ is primitive with probability $1-p^{-2}$,
in which case the determinant has moment $C_r$. Otherwise divide
that column by $p$ and repeat. It follows that
\[
 D_\infty(r)=\frac{1}{g_p}
 \left(N_0+N_1p^{-r}C_r+
 p^{-2r}\frac{(1-p^{-2})C_r}{1-p^{-2-r}}\right).
\]
For $p=5$ the values at $r=3$ and $4$ are $247551/324896$ and
$18565153/24404688$, respectively. This calculation uses the
division-tree framework of Jones--Rouse and Lombardo--Perucca
\cite{JR,LP}.

There is also a signed component. Let $\chi$ be the quadratic
character of the residual determinant and let $\Phi_r$ indicate
simultaneous divisibility modulo $p$. Then its inner product with
$\chi$ on the finite Galois group is
\[
 \langle\Phi_r,\chi\rangle=
 \frac{p-(p+1)p^{-r}+p^{-2r}}{g_p}>0 \quad (r>0).
\]
The signed counts for fixed-space dimensions $2,1,0$ are
$1,-(p+1),p$: the distinct-eigenvalue term sums the nontrivial
quadratic character over $\lambda\ne0,1$, giving $-p(p+1)$;
the unipotent term adds $p^2-1$. Thus taking only the trivial
character changes the Hecke-weighted prime sum even at the level
of its finite character expansion. No order of vanishing for the
original complex $L$-function follows from this density.

\section{A family with four marked directions}
Take the curves $E_t:y^2=x^3-28x+t^2-48$ and infinitely many distinct
$j$-invariants in the family of \cite{K3}. There the four points with
$x=-2,-4,0,-5$ specialize, at $t_0=817/56$, to $y$-coordinates
$817/56,817/56,719/56,751/56$ and have independent classes modulo
five. Independence persists on the congruence progression constructed
by the five good reductions in that paper. We refine that progression
at one additional prime to obtain full marked division fields.

At $t_0$ the discriminant is $-q_0/(2^8 7^4)$ with
$q_0=6352166814299$, a prime. For a short primality certificate,
$q_0-1=2\cdot398731\cdot7965479$; the two odd factors are prime by
trial division, $2^{q_0-1}=1\bmod q_0$, and
$\gcd(2^{(q_0-1)/h}-1,q_0)=1$ for each prime $h\mid q_0-1$.
These checks use repeated squaring. Require the parameter on the
elliptic curve $B$ of \cite{K3} to be congruent to $2T$ modulo $q_0^2$
as well. Multiples $(2+kM)T$ do so after multiplying the formerly
chosen $M$ by $q_0\#B(\F_{q_0})$, by the formal group at this good prime.
The inverse rational parameters at $2T=(100,224)$ are integral at
$q_0$. Thus $v_{q_0}(\Delta)=1$ and $c_4$ is a unit on the new
progression. Inertia at $q_0$ supplies a transvection on $E_t[5]$
by the Tate-curve description (over an unramified extension if needed).
At $13$ the unchanged trace is $-2$, with irreducible residual
characteristic polynomial, so the residual representation is
irreducible. An irreducible subgroup containing a transvection contains
$\operatorname{SL}_2(\F_5)$: some conjugate has distinct fixed line,
and in the basis of the two fixed lines the two transvections generate
the upper and lower unipotents. Cyclotomic determinant gives all of
$\GL_2(\F_5)$. The usual lifting lemma of Serre for $p\ge5$ then
gives full $5$-adic image \cite{Serre}. Equivalently apply the lemma
to the closure of the commutator subgroup in $\operatorname{SL}_2$
and then use the full cyclotomic determinant.

For these four marked points the field degree over $\Q$ at level
$5^n$ is $|\GL_2(\Z/5^n)|5^{8n}=480\cdot5^{12n-4}$.
Their simultaneous all-level divisibility density is $D_\infty(4)$.
If the full rational rank is larger than four, its all-point density
is smaller. This explains the distinction between marked reductions
and the canonical fields used to recover $\Sha[p^n]$.

\begin{thebibliography}{9}
\bibitem{JR} R. Jones and J. Rouse, \emph{Galois theory of iterated
endomorphisms}, \href{https://doi.org/10.1112/plms/pdp051}{Proc. London Math. Soc.}
\bibitem{LP} D. Lombardo and A. Perucca, \emph{Reductions of points
on algebraic groups}, \href{https://arxiv.org/abs/1612.02847}{arXiv:1612.02847}.
\bibitem{Serre} J.-P. Serre, \emph{Abelian $\ell$-adic representations and
elliptic curves}, Chap.~IV, \S3 (lifting lemma).
\bibitem{K3} \emph{Ranks of two elliptic K3 surfaces obtained by quadratic
base change}, \href{https://github.com/cosmosgeek0/BSD/blob/main/quadratic-base-change-k3.pdf}{companion paper}.
\end{thebibliography}
\end{document}
